\section{November 30 Exit Criteria}
\label{sec:exit-criteria}

Both architectures must pass. Each line is a specific, checkable claim.

\subsection{Greenfield exit gate}

\begin{itemize}[itemsep=0.1em, label={$\square$}]
\item Clicker~4 + LAN8651 functioning.
\item LAN8660 functioning.
\item LAN8661 functioning.
\item All three coexist on a T1S segment.
\item PLCA characterized.
\item LAN8660 controls a real motor-driver/actuator.
\item LAN8660 provides observable peripheral state.
\item LAN8661 controls a real physical lighting load.
\item Both endpoints operate simultaneously.
\item Latency measured.
\item Jitter measured.
\item Packet loss measured.
\item Node disconnect/reconnect demonstrated.
\item Power-cycle recovery demonstrated.
\item Current/power measurements recorded.
\end{itemize}

\subsection{Brownfield exit gate}

\begin{itemize}[itemsep=0.1em, label={$\square$}]
\item Controlled CAN network operational.
\item Alfred Gateway attached without replacing the network.
\item Passive mode hardware behavior verified.
\item Zero Gateway TX while passive verified independently.
\item CAN traffic captured.
\item Traffic timestamped.
\item Traffic normalized/mirrored onto T1S.
\item Clicker receives mirrored state.
\item Authorized bidirectional mode explicitly enabled.
\item Clicker command traverses: Clicker $\rightarrow$ T1S $\rightarrow$ Gateway $\rightarrow$ CAN $\rightarrow$ physical device.
\item Physical response occurs.
\item Response mirrored back.
\item Total round-trip latency measured.
\item CAN-FD validated if possible.
\end{itemize}

LIN remains optional if schedule pressure exists (Section~\ref{sec:feature-cuts}).

\begin{decision}[title={How this gate is scored}]
Every unchecked item is a named gap in the November~30 report, walked explicitly against Section~\ref{sec:measurement} data at the exit review (Section~\ref{sec:timeline}) -- not against demo-day impressions.
\end{decision}
